\documentclass[main]{subfiles}
\usepackage[sols]{workbook}
\numbering{ws2-1}

\begin{document}

\ifSubfilesClassLoaded{%
  \chaptermark{\chaptertwo}
}{}

\section{Taylor Approximations I} \label{ws2-1}


\begin{enumerate}

\begin{comment}
     \item 
  \begin{problem}
    Approximate $\sqrt{102}$ by hand.  
  \end{problem}

\end{comment}


%\begin{comment}
  \item


  \begin{problem}[1in]
    Approximate $\sqrt{102}$ by hand.  
  \end{problem}

  \begin{solution}
    It is reasonable to make the approximation $\sqrt{102} \approx \sqrt{100} = 10$, but let's improve our estimate further.  The closest integer to 102 for which we can
    easily compute the square root is 100.  So here is the strategy we will use:  Let $f(x) = \sqrt{x}$
    and consider its tangent line approximation at the point $(100, 10)$.  The tangent line at this
    point is given by $L(x) = f(100) + f'(100)(x-100)$, so we need to compute the derivative of $f$.
    By the power rule, the derivative of $f(x) = x^{1/2}$ is
    $$
      f'(x) = \frac{1}{2} x^{-1/2} \; = \; \frac{1}{2 \sqrt{x}} \qquad (x > 0).
    $$
    Then $f'(100) = \frac{1}{20}$, and therefore the line tangent to the graph of $f$ at the point
    $(100,10)$ is given by $L(x) = 10 + \frac{1}{20}(x - 100)$.  

    For $x$ ``near" 100, we expect $L(x)$ to be a reasonably good approximation for $f(x)$, which
    means we can approximate $\sqrt{102} = f(102)$ using $L(102)$.  But $L(102)$ is easy to calculate
    by hand: $L(102) =  10 + \frac{1}{20}(102 - 100) = 10.1$.  So our tangent line approximation gives
    $\sqrt{102} \approx 10.1$.  This is pretty close to the actual value of $\sqrt{102}$ which, rounded
    to the nearest millionth is $10.099505$.    

    The graph of $f(x) = \sqrt{x}$ is concave downward because
    $$
      f''(x) = - \frac{1}{4} x^{-3/2} \; = \; \frac{-1}{4 x^{3/2}}
    $$
    is negative for all $x > 0$.  Because the graph of $f$ is concave downward, 
    the graph of the tangent line referenced in the solution of the preceding
    part of the
    problem lies \textbf{above} the graph of $f$ (except of course at the point of 
    tangency).  Thus, our approximation of $\sqrt{102}$ in the previous 
    part of the problem
    is an {\bf overestimate}.  
 
Now let's try to get a 2nd degree polynomial approximation. 

     First, we want to ensure that the values of $P_{2}$ and $f$ agree when $x = 100$. 
     Since $P_{2}(100) = a_0$ and $f(100) = 10$, this means we need \fbox{$a_0 = 10$}.

     Next, we want to ensure that the derivatives of $P_{2}$ and $f$ agree when $x = 100$.
     Since $P_{2}'(x) = a_{1} + 2a_{2}(x - 100)$, we see that $P_{2}'(100) = a_{1}$.
     Now since $f'(x) = \frac{1}{2 \sqrt{x}}$, we have $f'(100) = \frac{1}{20}$.  Thus,
     in order for $P_{2}'(100) = f'(100)$, we must require that \fbox{$a_{1} = \frac{1}{20}$}.

     Finally, we want to ensure that the second derivatives of $P_{2}$ and $f$ agree when $x = 100$.
     Since $P_{2}''(x) = 2a_{2}$, we see that $P_{2}''(100) = 2a_{2}$.
     Now since $f''(x) = -\frac{1}{4 x^{3/2}}$, we have $f''(100) = -\frac{1}{4000}$.  Thus,
     in order for $P_{2}''(100) = f''(100)$, we must require that \fbox{$a_{2} = -\frac{1}{8000}$}.

     It follows that
     $$
       P_{2}(x) = 10 + \frac{1}{20}(x-100) - \frac{1}{8000}(x-100)^2.
     $$
     For $x$ ``near" 100, we expect $P_{2}(x)$ to be a reasonably good approximation for $f(x)$, which
     means we can approximate $\sqrt{102} = f(102)$ using 
     $$
       P_{2}(102) = 10 + \frac{1}{20}(102-100) - \frac{1}{8000}(102-100)^2 \; = \; 10 + \frac{1}{10} - \frac{1}{2000}
                  \; = \; 10.0995.
     $$
     This turns out to be accurate to five decimal places!

  \end{solution}



%\end{comment}
%%%%%%%%%%%%%%%%%%%%%%%%%%%%%%%%%%%%%%%%%%%%%%%%%%%%%%%%%%%%%%%%%%%%%%%%%%%%%%%%%%%%%  

\pspace{\newpage}

\item[]
\begin{framed}
\textbf {\underline{Taylor polynomials:}}  Let $f$ be a function whose first $n$
derivatives exist at $c$.  The polynomial
$$
  P_{n}(x) = a_{0} + a_{1}(x-c) + a_{2}(x-c)^{2} + \cdots + a_{n}(x-c)^{n}
$$
for which the coefficients $a_0, a_1, \dots, a_n$ are obtained by requiring
$$
P_{n}(c) = f(c), \qquad P_{n}'(c) = f'(c), \qquad P_{n}''(c) = f''(c), \qquad \dots \qquad
P_{n}^{(n)}(c) = f^{(n)}(c)
$$
is called the \textbf {n-th degree Taylor polynomial for $f$, centered at $c$}.  
\end{framed}


\item 
  \note{In this example, $P_{3}(x)$ and $P_{4}(x)$ are actually the same.  Hence the footnote
  in the worksheet.  In Part (a), I want to emphasize the process of matching derivatives.  I'm
  not worried about giving a general formula for the coefficients, or mentioning factorials---they'll
  see all that next time.  In Part (b), I want to point out that if they have an $n$th-degree
  Taylor polynomial centered at $c$, then finding the $(n+1)$st-degree Taylor polynomial centered
  at $c$ only amounts to finding the coefficient of $(x-c)^{n+1}$; none of the lower coefficients
  change.}

\begin{enumerate}
  \item \label{sine}
  \begin{problem}[3in]
     Find the 3rd-degree Taylor polynomial for $f(x) = \sin(x)$, centered at $0$.
     Then, use your result to estimate $\sin(0.3)$.
  \end{problem}

  \begin{solution}
     We need to compute the coefficients of the polynomial
     $$
       P_{3}(x) = a_0 + a_1 x + a_2 x^2 + a_3 x^3,
     $$
     which is centered at $c = 0$.  To do so, we must require that
     $$
     P_{3}(0) = f(0), \qquad P_{3}'(0) = f'(0), \qquad P_{3}''(0) = f''(0), \qquad \mbox{and} \qquad
     P_{3}'''(0) = f'''(0),
     $$
     meaning that we'll need to compute the first three derivatives of $P_{3}$ and of $f$.

     First, we'll require $P_{3}(0) = f(0)$.
     Since $P_{3}(0) = a_0$ and $f(0) = \sin(0) = 0$, we have \fbox{$a_0 = 0$}.

     Next, we'll match the first derivatives, requiring $P_{3}'(0) = f'(0)$.  
     Now $P_{3}'(x) = a_1 + 2a_2 x + 3a_3 x^2$,
     so $P_{3}'(0) = a_1$.  And since $f'(x) = \cos(x)$, we have $f'(0) = 1$.  Setting
     $P_{3}'(0) = f'(0)$ yields \fbox{$a_1 = 1$}.

     Next, we'll match the second derivatives, requiring $P_{3}''(0) = f''(0)$.  
     Now $P_{3}''(x) = 2a_2  + 6 a_3 x$,
     so $P_{3}''(0) = 2a_2$.  And since $f''(x) = -\sin(x)$, we have $f''(0) = 0$.  Setting
     $P_{3}''(0) = f''(0)$ yields \fbox{$a_2 = 0$}.

     Finally, we'll match the third derivatives, requiring $P_{3}'''(0) = f'''(0)$.  
     Now $P_{3}'''(x) = 6 a_3$,
     so $P_{3}'''(0) = 6 a_3$.  And since $f'''(x) = -\cos(x)$, we have $f'(0) = -1$.  Setting
     $P_{3}'''(0) = f'''(0)$ yields \fbox{$a_3 = - \frac{1}{6}$}.

     Therefore, \fbox{$\displaystyle P_{3}(x) = x - \frac{1}{6} x^{3}$}.

     Because $P_{3}(x)$ was designed by trying to mimic the behavior of $f(x)$ near
     the center $c = 0$, we hope that $P_{3}(x)$ is a good approximation for $f(x)$ for
     $x$ near $0$.  So we'll approximate $f(0.3) = \sin(0.3)$ using
     $$
       P_{3}(0.3) = 0.3 - \frac{(0.3)^3}{6} \; = \; 0.2955.
     $$
     It happens that $\sin(0.3) = 0.2955202...$, so you may be surprised to find that
     the above approximation is actually accurate to four decimal places.

     The figure below shows an overlay of the graphs of $f$ and $P_{3}$ so that
     you can see how closely they agree near the center $c = 0$.
     \begin{center}
        \includegraphics[width=2in]{ws/images/sin-cubic.png}
     \end{center}

  \end{solution}
  \item

  \begin{problem}[0.5in]
     Find the 4th-degree Taylor polynomial for $f(x) = \sin(x)$, centered at $0$.
     Then, use your results to estimate $\sin(0.3)$.
  \end{problem}

  \begin{solution}
     We've done most of the work already.  If
     $$
       P_{4}(x) = a_0 + a_1 x + a_2 x^2 + a_3 x^3 + a_4 x^4,
     $$
     the coefficients $a_0$, $a_1$, $a_2$, and $a_3$ were computed in the previous
     part by matching up zeroth, first, second and third derivatives.  All that
     remains is to match fourth derivatives, requiring $P_{4}^{(4)}(0) = f^{(4)}(0)$.  
     Now $P_{4}^{(4)}(x) = 24 a_4$,
     so $P_{4}^{(4)}(0) = 24 a_4$.  And since $f^{(4)}(x) = \sin(x)$, we have $f^{(4)}(0) = 0$.  
     Setting $P_{4}^{(4)}(0) = f^{(4)}(0)$ yields \fbox{$a_4 = 0$}.

     It follows that $P_{4}(x)$ is the {\bf same} polynomial as $P_{3}(x)$.  
     In the footnote, we prepped you to expect this sort of thing.
  \end{solution}

\end{enumerate}



  
%%%%%%%%%%%%%%%%%%%%%%%%%%%%%%%%%%%%%%%%%%%%%%%%%%%%%%%%%%%%%%%%%%%%%%%%%%%%%%%%%%%%%  



\end{enumerate}


\end{document}